% Responsable: José
\subsection{Plataforma de una cooperativa}

\subsubsection{Caracterización}

El enunciado únicamente menciona una plataforma de una cooperativa, sin especificar sus funciones ni condiciones de desarrollo. Para realizar la comparación, se plantea un caso hipotético de una cooperativa de ahorro y crédito. Las características y restricciones descritas son supuestos de trabajo, no datos obtenidos de una organización real.

\paragraph{Descripción de la plataforma}

Se supone que la cooperativa necesita una plataforma web para que sus socios consulten saldos y movimientos, realicen transferencias entre cuentas de la misma institución y presenten solicitudes de crédito. Los empleados podrán revisar las solicitudes y registrar su seguimiento; la aprobación del crédito continuará a cargo del personal autorizado.

La plataforma se conectará con un sistema existente que mantiene las cuentas y registra las operaciones. No sustituirá ese sistema ni calculará directamente los saldos. Quedan fuera del alcance inicial las transferencias a otras instituciones, la aplicación móvil y la aprobación automática de créditos.

\paragraph{Cliente y actores del proceso}

El cliente será la cooperativa, representada por su gerencia. Se consideran los siguientes actores:

\begin{itemize}
	\item \textbf{Gerencia:} acuerda el alcance, las prioridades y la aceptación del producto.
	
	\item \textbf{Responsables de operaciones y crédito:} explican los procedimientos, las reglas del negocio y las excepciones que debe contemplar la plataforma.
	
	\item \textbf{Socios:} aportan necesidades de uso y participan en la revisión de consultas, transferencias y solicitudes.
	
	\item \textbf{Responsable de tecnología:} proporciona información sobre el sistema existente y sus interfaces.
	
	\item \textbf{Responsable de seguridad y cumplimiento:} revisa las necesidades de protección de datos, permisos y registro de operaciones.
	
	\item \textbf{Equipo de desarrollo:} obtiene, documenta y analiza los requisitos, identifica conflictos y prepara los incrementos para su validación.
\end{itemize}

\paragraph{C1. Estabilidad de los requisitos: media}

Se supone que las operaciones principales y las reglas para registrar movimientos ya están definidas en el sistema existente. Sin embargo, los detalles de las pantallas, el seguimiento de solicitudes y las notificaciones necesitan aclararse con empleados y socios. Por ello, existe una base estable, pero parte de los requisitos puede ajustarse durante las revisiones.

\paragraph{C2. Criticidad del sistema: alta}

La plataforma permitirá acceder a información financiera e iniciar transferencias. Una falla podría exponer datos, duplicar una operación o mostrar información incorrecta. Por esta razón, el escenario requiere comprobar los permisos de acceso, la integridad de las operaciones y la correspondencia de lo mostrado con los registros del sistema existente.

\paragraph{C3. Disponibilidad del cliente: media}

Se supone que la cooperativa designará a un representante de operaciones para una reunión semanal y para responder dudas entre reuniones. Los responsables de crédito, seguridad y tecnología participarán en revisiones programadas. Los socios estarán disponibles para sesiones puntuales de validación, pero no para colaborar diariamente.

Esta disponibilidad permite recibir retroalimentación periódica, aunque algunas decisiones tendrán que esperar la intervención del responsable correspondiente.

\paragraph{C4. Tamaño y complejidad: mediano}

Se plantea un equipo de cinco personas para desarrollar una plataforma con tres grupos principales de funciones: consultas, transferencias internas y solicitudes de crédito. Participan usuarios con distintos permisos y existe una integración con el sistema de cuentas.

El alcance se considera mediano por estar limitado a esas funciones. No obstante, la integración y el manejo de operaciones financieras añaden complejidad técnica y requieren pruebas específicas.

\paragraph{C5. Regulación y contrato: alta}

Como supuesto del caso, la cooperativa exigirá revisiones de seguridad y cumplimiento antes de autorizar el uso de funciones que manejen datos u operaciones financieras. Las obligaciones normativas concretas deberán identificarse y verificarse antes de establecer la línea base de requisitos.

También se supone un contrato con presupuesto y fecha de entrega acordados, organizado por entregables. Los cambios que afecten alcance, costo o plazo deberán registrarse y aprobarse. Estas restricciones exigen controlar las modificaciones, aunque permiten revisar los requisitos dentro de los límites acordados.

\paragraph{C6. Documentación exigida: alta}

Se supone que la cooperativa solicitará una especificación de requisitos, reglas de negocio, criterios de aceptación, descripción de interfaces y registros de cambios. También exigirá trazabilidad entre los requisitos y sus pruebas, junto con evidencias de validación y aprobación de los entregables.

Esta documentación permitirá comprobar qué se acordó, por qué se modificó y con qué resultados se aceptó cada función.

\begin{table}[H]
	\centering
	\footnotesize
	\caption{Caracterización del escenario hipotético de la cooperativa}
	\label{tab:escenario-cooperativa}
	\begin{tabularx}{\textwidth}{@{}Yp{0.22\textwidth}@{}}
		\toprule
		\textbf{Criterio} & \textbf{Valoración} \\
		\midrule
		C1. Estabilidad de los requisitos & Media \\
		C2. Criticidad del sistema & Alta \\
		C3. Disponibilidad del cliente & Media \\
		C4. Tamaño y complejidad & Mediano \\
		C5. Regulación y contrato & Alta \\
		C6. Documentación exigida & Alta \\
		\bottomrule
	\end{tabularx}
\end{table}

Estas valoraciones describen el escenario hipotético. Se utilizarán como base para comparar los modelos en cascada, iterativo y ágil, sin asumir de antemano que alguno es el más adecuado.

\subsubsection{Modelo elegido y justificación}

\textbf{Modelo elegido: iterativo con desarrollo incremental y revisiones documentadas.}

Se propone construir versiones sucesivas, revisar los requisitos con la cooperativa y ampliar las funciones de forma controlada. Se distingue entre iterar, que permite revisar y mejorar lo trabajado, e incrementar, que añade funcionalidad. Ambas prácticas se combinan en la propuesta para este caso \cite{sommerville2016}.

La elección se apoya en las siguientes razones del escenario:
\begin{enumerate}
    \item \textbf{Hay requisitos que todavía necesitan aclararse.} Las operaciones principales tienen una base definida, pero las pantallas, las notificaciones y el seguimiento de solicitudes pueden cambiar tras las revisiones. Los ciclos permiten precisar esos detalles antes de extenderlos a toda la plataforma.
    \item \textbf{La integración puede comprobarse progresivamente.} Se plantea comenzar con consultas de saldos y movimientos para revisar la conexión con el sistema existente, los permisos y la correspondencia de los datos. Después se incorporarían las transferencias internas y las solicitudes de crédito. Cada incremento necesitará pruebas y aprobación antes de ponerse en operación.
    \item \textbf{La disponibilidad del cliente permite revisiones programadas.} La reunión semanal con operaciones puede utilizarse para aclarar dudas y revisar avances. Las decisiones de crédito, tecnología y seguridad se agendarán antes del cierre de cada ciclo. Se supone que estos responsables pueden cumplir esos compromisos; si no ocurre, deberá revisarse la planificación.
    \item \textbf{La documentación puede mantenerse por incremento.} Antes de aceptar cada entrega se actualizarán los requisitos, las interfaces, los criterios de aceptación, la trazabilidad y los resultados de pruebas. Los cambios con impacto en alcance, costo o plazo se someterán a la aprobación prevista en el contrato supuesto.
\end{enumerate}

La matriz de la Tabla~\ref{tab:matriz-cooperativa} respalda esta elección con 17 puntos. La ventaja sobre cascada es reducida y depende principalmente de poder aclarar requisitos e integrar el sistema de cuentas por partes sin abandonar las revisiones formales.

\textbf{Alternativas descartadas para este escenario:}

\textbf{Cascada.} Es una alternativa viable por el control de cambios, la documentación y las aprobaciones que exige la cooperativa. No se elige como organización principal porque todavía quedan requisitos de uso por precisar y conviene comprobar pronto la integración con el sistema existente. Esperar a una etapa posterior para validar estos aspectos aumentaría el retrabajo. Su valoración de 16 puntos muestra que sigue siendo una opción cercana, especialmente si se aclararan los requisitos pendientes antes de construir.

\textbf{Ágil.} También permite revisar requisitos y entregar incrementos, por lo que responde bien a la estabilidad media y al tamaño del proyecto. Se considera menos conveniente en las condiciones asumidas porque socios y especialistas no estarán disponibles de forma continua y deben coordinarse aprobaciones documentales y contractuales. El enfoque ágil puede incorporar esos controles; su descarte no se basa en que carezca de documentación, sino en que la propuesta iterativa planificada se ajusta mejor a la participación programada y a los entregables formales del caso \cite{schon2017,turk2014}.

\subsubsection{Riesgos del modelo elegido y mitigación}

Los siguientes riesgos se plantean a partir de los supuestos de la cooperativa y de la propuesta de desarrollar la plataforma mediante ciclos e incrementos. No representan problemas observados en una organización real.

\begin{enumerate}
    \item \textbf{Retraso de las revisiones.} Si operaciones, crédito o seguridad no revisan un incremento a tiempo, el siguiente ciclo podría comenzar con requisitos pendientes y generar retrabajo. \textbf{Mitigación:} acordar las fechas de revisión al inicio de cada ciclo, enviar previamente los puntos por resolver y registrar las decisiones. No iniciar funciones que dependan de una aprobación pendiente; continuar con tareas que ya tengan requisitos acordados.

    \item \textbf{Crecimiento del alcance entre ciclos.} Las revisiones de pantallas y solicitudes pueden generar nuevas funciones que superen el presupuesto y la fecha pactados. \textbf{Mitigación:} distinguir las aclaraciones de los requisitos de las nuevas solicitudes, estimar su impacto y obtener aprobación antes de incorporarlas. Mantener fuera del alcance inicial las transferencias externas, la aplicación móvil y la aprobación automática de créditos, salvo cambio contractual autorizado.

    \item \textbf{Errores de integración y regresión.} Un incremento de transferencias podría alterar las consultas ya validadas o registrar una operación de forma duplicada al comunicarse con el sistema de cuentas. \textbf{Mitigación:} comprobar la interfaz desde el primer incremento en un entorno de pruebas y repetir las pruebas de las funciones anteriores en cada ciclo. Incluir casos de interrupción, reintento y respuesta fallida; verificar que un reintento no duplique la operación y que saldos y movimientos coincidan con el sistema existente.

    \item \textbf{Controles de seguridad incompletos en entregas parciales.} Al priorizar funciones visibles, el equipo podría posponer permisos o registros necesarios para proteger las consultas y transferencias. \textbf{Mitigación:} incluir requisitos y criterios de aceptación de acceso y registro de operaciones desde el primer ciclo. El responsable de seguridad deberá revisar las evidencias antes de autorizar cada entrega. Una versión mostrada para recibir comentarios no se considerará automáticamente aprobada para operar con datos o dinero reales.

    \item \textbf{Documentación desactualizada.} Los ajustes de cada iteración pueden dejar diferencias entre lo implementado, los requisitos aprobados y las pruebas, dificultando la aceptación del entregable. \textbf{Mitigación:} actualizar la especificación, las interfaces, la trazabilidad y los resultados de pruebas dentro del mismo ciclo. Comprobar su correspondencia con la versión entregada y registrar la aprobación antes de cerrar el incremento.
\end{enumerate}
